\section{Training und mobile Umsetzung}\label{sec:training}

\subsection{Trainingsrezept}\label{sec:rezept}

Ein Python-Skript trainiert die vier Modelle der drei Ansätze mit
Ultralytics 8.3.155 \cite{jocher2024ultralytics}. Es läuft auf
NVIDIA-GPUs mit CUDA, auf dem Mac mit Apple MPS oder auf der CPU.
Als Ausgangspunkt dienen
vortrainierte Gewichte von YOLO11n für die Detektion, YOLO11n-obb für
orientierte Boxen und YOLO11n-cls für die Klassifikation.

Ein Lauf umfasst normalerweise 20 Epochen mit einer Bildgrösse von
640 Pixeln. Auf dem Mac werden 16 Bilder pro Batch verarbeitet, auf der
Cloud-GPU 64. Ultralytics erhöht die Lernrate während der ersten drei
Epochen schrittweise. Bei den Detektoren wird die Mosaic-Augmentierung
in den letzten zehn Epochen abgeschaltet. Die Voreinstellung für einen
vorzeitigen Abbruch liegt bei 100 Epochen und greift bei diesen Läufen
nicht \cite{ultralytics2026cfg}.

Horizontale und vertikale Spiegelungen sind deaktiviert, weil eine
gespiegelte Karte eine andere Karte ergibt. Die Klassennamen stammen
aus der gemeinsamen Kartenliste. Bei den Klassifikatoren ordnet
Ultralytics die Klassen intern alphabetisch. Deshalb lesen die Apps
die zum Modell gehörenden Klassennamen, statt sich auf die
Klassenindizes zu verlassen.

Trainiert wird auf synthetischen Bildern. Nach jeder Epoche wird auf
realen Fotos validiert. Die Gewichte der besten Epoche werden als
\code{best.pt} gespeichert. Weil diese Epoche anhand desselben
Validierungssets gewählt wird, auf dem die Ergebnisse berichtet
werden, können die Kennzahlen etwas zu hoch ausfallen
(Abschnitt~\ref{sec:grenzen}).

\subsection{Reproduzierbare Trainingspipeline}\label{sec:pipeline}

Die ersten Modelle wurden auf einem MacBook trainiert. Für grössere
Datensätze und den Vergleich aller Varianten kamen gemietete
Cloud-GPUs zum Einsatz. Ein GitHub-Actions-Workflow startet dafür
einen RunPod-Pod. Dieser lädt den festgelegten Commit, erzeugt den
Datensatz aus Seed und Bildanzahl und trainiert die Varianten
nacheinander. Während des Laufs werden Manifest, Log, Gewichte und
Metriken in Azure Blob Storage gespeichert. Ein lokales Werkzeug
holt die Ergebnisse zurück und exportiert das gewählte Modell für
beide Apps.

Die Trainingsbilder werden für jeden Lauf neu erzeugt. Grundlage sind
der versionierte Code, die Kartenscans und die festgelegten
Generatorparameter (Abschnitt~\ref{sec:repro}). Das Manifest hält
Commit, Seed, Bildanzahl und Konfiguration fest. Ein Wächter prüft
halbstündlich, ob die Läufe noch Fortschritte melden, und beendet
verwaiste Pods. Die Pipeline ist in
\code{src/training/pipeline.md} dokumentiert.

\subsection{Export und mobile Inferenz}\label{sec:app}

Aus \code{best.pt} entstehen zwei Modelle mit denselben trainierten
Gewichten. Für die eingereichten Apps wird Variante~C verwendet.
Ihr Core-ML-Modell enthält die Non-Maximum Suppression, sodass Vision
fertige Boxen liefert. Ultralytics speichert dieses Modell in halber Genauigkeit
(FP16) \cite{ultralytics2026coreml}. Für Android entsteht ein LiteRT-Modell in voller Genauigkeit
(FP32) ohne diesen Schritt. Die App wertet dessen Ausgabe selbst
aus. Bei der Berechnung auf einer Android-GPU kann LiteRT
trotzdem mit FP16 rechnen \cite{google2026litertgpu}.

Eine zusätzliche Quantisierung auf 8 Bit wurde nicht vorgenommen.
Das Core-ML-Modell ist 5.3\,MB gross. Das LiteRT-Modell benötigt aufgrund
der vollen Genauigkeit 10.7\,MB. Um die Geschwindigkeit in der App
einzuschätzen, wurde die Laufzeit pro Bild auf vier verfügbaren
Smartphones gemessen (Tabelle~\ref{tab:geraete}). Auf den beiden iPhones
dauert die Erkennung 15 beziehungsweise 30\,ms. Das Pixel~7a erreicht mit
50\,ms genau den festgelegten Zielwert. Auf dem Nexus~5X benötigt die
Verarbeitung dagegen rund drei Sekunden pro Bild. Damit lässt sich die
App auf diesem Gerät nicht sinnvoll bedienen.

Die Werte sind gerundete Einzelmessungen und dienen nur als grobe
Einordnung. Sie erlauben keinen allgemeinen Vergleich zwischen iOS und
Android. Arbeitsspeicher und Energieverbrauch wurden nicht gemessen. Für
schwächere Android-Geräte bleibt eine Quantisierung auf 8 Bit eine
mögliche Erweiterung (Abschnitt~\ref{sec:ausblick}).

\begin{table}[H]
\centering\small
\begin{tabular}{@{}lclr@{}}
\toprule
\textbf{Gerät} & \textbf{Jahr} & \textbf{Laufzeitumgebung} & \textbf{Zeit pro Bild} \\
\midrule
iPhone 17 Pro Max & 2025 & Core ML & 15\,ms \\
iPhone 14 Pro & 2022 & Core ML & 30\,ms \\
Google Pixel 7a & 2022 & LiteRT & 50\,ms \\
Nexus 5X & 2015 & LiteRT & 3000\,ms \\
\bottomrule
\end{tabular}
\caption{Gerundete Laufzeit der Variante~C pro Bild. Die Werte stammen aus
einer Einzelmessung in der jeweiligen App.}
\label{tab:geraete}
\end{table}

Die Klassennamen sind Teil des Exports. Die Datei \code{models.json}
nennt für jedes gebündelte Modell den Trainingslauf und den Commit.

Die App wurde mit SwiftUI für iOS und mit Kotlin und Jetpack Compose
für Android umgesetzt. Beide verwenden im Kamerabild das grösste
mittige Quadrat und skalieren es auf 640 Pixel. Damit entspricht
der Bildausschnitt dem Format der Trainingsbilder. Für den
Variantenvergleich konnte die iOS-Implementierung auch A und B
verarbeiten. Bei B entzerrt sie den von der ersten Stufe gefundenen
Ausschnitt und klassifiziert ihn danach. Die eingereichten Apps
enthalten Variante~C.

\subsection{Von der Erkennung zur gezählten Karte}\label{sec:zaehlung}

Aus dem geplanten Proof of Concept ist eine fertig implementierte
Zähl-App für iOS und Android geworden. Beide Versionen wurden am
14.\,September als Version 1.0.0 im App Store und bei Google Play
eingereicht. Sie sind als Gratis-Demo mit einmaligem Kauf angelegt.
Die Kaufabwicklung, Store-Texte und Screenshots gingen über den
geplanten Proof of Concept und den Untersuchungsrahmen dieser Arbeit
hinaus. Ziel ist, die App interessierten Personen am Tag der
Präsentation zugänglich zu machen. Die guten Erkennungsresultate
sprechen dafür, die App danach weiterzuentwickeln und zu pflegen.

Eine Erkennung wird noch nicht sofort gezählt. Die App verwirft
Vorhersagen mit einer Konfidenz unter 0.6. Danach prüft sie, ob
dieselbe Karte in mehreren Frames erscheint. Standardmässig muss
sie in drei aufeinanderfolgenden Frames die beste Erkennung sein
(\enquote{Serie}). Als Alternative reichen drei von fünf Frames,
solange keine andere noch ungezählte Karte mehr als einmal
erscheint (\enquote{Mehrheit}). Die benötigte Anzahl lässt sich
von 1 bis 10 einstellen. Schwelle und Standardeinstellung der
Stabilitätsregel wurden zunächst ohne Messung gewählt. Ihre
Wirkung wurde später an einer aufgezeichneten Session geprüft
(Abschnitt~\ref{sec:betriebspunkt}).

Vor der Zählung werden das Kartenblatt und die Spielart gewählt. Danach
öffnet sich die Kamera. Die App
berücksichtigt nur die 36 Karten dieses Blatts und verhindert
Doppelzählungen. Der Stapel lässt sich durch Antippen korrigieren.
Die Punkte werden nach neun Spielarten aus drei Wertetabellen
berechnet. Ein Prüfskript kontrolliert diese Tabellen in der
Continuous Integration.
Für jede gezählte Karte gibt die App ein Signal. Ein verborgener
Entwicklermodus kann die Zählung als Video und als
Erkennungsprotokoll je Frame aufzeichnen. Diese Aufnahmen dienen
der Auswertung in Abschnitt~\ref{sec:betriebspunkt}.
